\renewcommand{\bibname}{Références}
\begin{thebibliography}{99}
\addcontentsline{toc}{chapter}{Références}

\bibitem{baac}
ONISR / ministère de l'Intérieur.
\textit{Bases de données annuelles des accidents corporels de la circulation routière,
années de 2005 à 2024}. Publication sur
\href{https://www.data.gouv.fr/datasets/bases-de-donnees-annuelles-des-accidents-corporels-de-la-circulation-routiere-annees-de-2005-a-2024}{data.gouv.fr}.
Le projet utilise les fichiers 2019-2024 fournis et identifiés par son manifeste.

\bibitem{doc-onisr}
ONISR. \textit{Description des bases de données annuelles des accidents corporels de la
circulation routière, années 2005 à 2024}. 21 octobre 2025.
Document fourni avec les données ; copie incluse dans les sources du rapport.
Voir notamment les pages imprimées 1-3, 7, 10-11 pour les limites de collecte,
les unités, les véhicules sans usager et les modalités de gravité.

\bibitem{discussion-lieux}
DSR ONISR. Réponse du 12 décembre 2024 à la discussion «~lieux-2023~».
\href{https://www.data.gouv.fr/datasets/bases-de-donnees-annuelles-des-accidents-corporels-de-la-circulation-routiere-annees-de-2005-a-2024/discussions}{Discussions de la page BAAC, data.gouv.fr}.
Consultée le 16 octobre 2026.

\bibitem{stata-bsample}
StataCorp. \textit{bsample: Sampling with replacement}.
\href{https://www.stata.com/manuals/rbsample.pdf}{Manuel officiel, options cluster et strata}.
Consulté le 16 octobre 2026.

\bibitem{nist-bootstrap}
NIST / SEMATECH. \textit{e-Handbook of Statistical Methods}, section 1.3.3.4,
\href{https://itl.nist.gov/div898/handbook/eda/section3/eda334.htm}{Bootstrap Plot}.
Consulté le 16 octobre 2026.

\bibitem{notebooks-projet}
Papa Birane Mbengue. \textit{Projet BAAC : notebooks 01 à 07 et résultats agrégés}.
Version fournie et vérifiée le 17 octobre 2026. Découverte, audit, préparation,
exploration, statistique, validation 2023 et test rétrospectif 2024.
Dépôt : \href{https://github.com/birane-m/road-accident-severity-analysis}{road-accident-severity-analysis}.
Les chiffres et graphiques du rapport proviennent des exports de cette version.

\bibitem{sklearn-logistique}
scikit-learn. \textit{Linear Models, Logistic regression}, documentation 1.9.1.
\href{https://scikit-learn.org/stable/modules/linear_model.html#logistic-regression}{Modèle binaire, objectif régularisé et choix des solveurs}.
Consultée le 17 octobre 2026.

\bibitem{sklearn-logloss}
scikit-learn. \textit{log\_loss}, documentation 1.9.1.
\href{https://scikit-learn.org/stable/modules/generated/sklearn.metrics.log_loss.html}{Définition et calcul de la perte logarithmique}.
Consultée le 17 octobre 2026.

\bibitem{sklearn-brier}
scikit-learn. \textit{brier\_score\_loss}, documentation 1.9.1.
\href{https://scikit-learn.org/stable/modules/generated/sklearn.metrics.brier_score_loss.html}{Définition du score de Brier}.
Consultée le 17 octobre 2026.

\bibitem{sklearn-ap}
scikit-learn. \textit{average\_precision\_score}, documentation 1.9.1.
\href{https://scikit-learn.org/stable/modules/generated/sklearn.metrics.average_precision_score.html}{Average precision et distinction avec l'aire trapézoïdale}.
Consultée le 17 octobre 2026.

\bibitem{sklearn-roc}
scikit-learn. \textit{roc\_auc\_score}, documentation 1.9.1.
\href{https://scikit-learn.org/stable/modules/generated/sklearn.metrics.roc_auc_score.html}{Aire sous la courbe ROC}.
Consultée le 17 octobre 2026.

\bibitem{sklearn-calibration}
scikit-learn. \textit{Probability calibration}, documentation 1.9.1.
\href{https://scikit-learn.org/stable/modules/calibration.html}{Calibration, diagrammes de fiabilité et limites des scores globaux}.
Consultée le 17 octobre 2026.

\end{thebibliography}
